% Options for packages loaded elsewhere
\PassOptionsToPackage{unicode}{hyperref}
\PassOptionsToPackage{hyphens}{url}
\PassOptionsToPackage{dvipsnames,svgnames,x11names}{xcolor}
\documentclass[
  11pt,
]{article}
\usepackage{xcolor}
\usepackage[margin=25mm]{geometry}
\usepackage{amsmath,amssymb}
\setcounter{secnumdepth}{-\maxdimen} % remove section numbering
\usepackage{iftex}
\ifPDFTeX
  \usepackage[T1]{fontenc}
  \usepackage[utf8]{inputenc}
  \usepackage{textcomp} % provide euro and other symbols
\else % if luatex or xetex
  \usepackage{unicode-math} % this also loads fontspec
  \defaultfontfeatures{Scale=MatchLowercase}
  \defaultfontfeatures[\rmfamily]{Ligatures=TeX,Scale=1}
\fi
\usepackage{lmodern}
\ifPDFTeX\else
  % xetex/luatex font selection
\fi
% Use upquote if available, for straight quotes in verbatim environments
\IfFileExists{upquote.sty}{\usepackage{upquote}}{}
\IfFileExists{microtype.sty}{% use microtype if available
  \usepackage[]{microtype}
  \UseMicrotypeSet[protrusion]{basicmath} % disable protrusion for tt fonts
}{}
\makeatletter
\@ifundefined{KOMAClassName}{% if non-KOMA class
  \IfFileExists{parskip.sty}{%
    \usepackage{parskip}
  }{% else
    \setlength{\parindent}{0pt}
    \setlength{\parskip}{6pt plus 2pt minus 1pt}}
}{% if KOMA class
  \KOMAoptions{parskip=half}}
\makeatother
\usepackage{listings}
\newcommand{\passthrough}[1]{#1}
\lstset{defaultdialect=[5.3]Lua}
\lstset{defaultdialect=[x86masm]Assembler}
\usepackage{longtable,booktabs,array}
\usepackage{calc} % for calculating minipage widths
% Correct order of tables after \paragraph or \subparagraph
\usepackage{etoolbox}
\makeatletter
\patchcmd\longtable{\par}{\if@noskipsec\mbox{}\fi\par}{}{}
\makeatother
% Allow footnotes in longtable head/foot
\IfFileExists{footnotehyper.sty}{\usepackage{footnotehyper}}{\usepackage{footnote}}
\makesavenoteenv{longtable}
\setlength{\emergencystretch}{3em} % prevent overfull lines
\providecommand{\tightlist}{%
  \setlength{\itemsep}{0pt}\setlength{\parskip}{0pt}}
\emergencystretch=3em
\AtBeginDocument{\Urlmuskip=0mu plus 2mu\urlstyle{tt}}
\lstset{breaklines=true,basicstyle=\small\ttfamily,columns=fullflexible,keepspaces=true}
\setlength{\tabcolsep}{4pt}
\usepackage{bookmark}
\IfFileExists{xurl.sty}{\usepackage{xurl}}{} % add URL line breaks if available
\urlstyle{same}
\hypersetup{
  pdftitle={A Lexically Filtered Meron-Cluster Benchmark: Exact Ground Truth, Mechanized Cancellation Lemmas, and a Pilot That Exposes Answer Leakage},
  pdfauthor={Jeremy H. Carroll},
  colorlinks=true,
  linkcolor={Maroon},
  filecolor={Maroon},
  citecolor={Blue},
  urlcolor={Blue},
  pdfcreator={LaTeX via pandoc}}

\title{A Lexically Filtered Meron-Cluster Benchmark: Exact Ground Truth,
Mechanized Cancellation Lemmas, and a Pilot That Exposes Answer Leakage}
\author{Jeremy H. Carroll}
\date{2026-10-03 · Version 0.1.0 · Creative Commons Attribution 4.0
International (CC-BY-4.0)}

\begin{document}
\makeatletter
\begin{center}
{\LARGE \@title\par}
\vspace{0.8em}
{\normalsize \@author\par}
\vspace{0.5em}
{\small \@date\par}
\end{center}
\makeatother

\subsection{Abstract}\label{abstract}

This paper packages one small, exactly solvable instance of the
meron-cluster mechanism of Chandrasekharan and Wiese and reports a pilot
language-model run on it that cannot support a discovery claim. The
instance is a four-site periodic chain with 256 breakup graphs and 1,296
configurations, of which 192 are negative. An exact oracle built on
SQLite and rational arithmetic finds that the 57 graphs containing a
meron cluster cancel exactly and that the other 199 give the signed sum
\(Z = 912\); an independent transfer-matrix trace also gives 912. The
per-orbit cancellation lemmas are checked in Lean 4. A substring filter
and five prompt builders pose the instance as a graph problem without
twelve domain words. In the retained pilot (five models accessed through
OpenRouter, one run per model and arm), no unassisted prompt produced a
passing answer: 0 of 9 cells. Every pass came after the prompt contained
the answer. The scorer's feedback states ``Z must equal 912'' and the
split into 57 cancelling and 199 surviving subgraphs, and the fifth arm
also dictates the sign formula and the cancellation step. Even the
unassisted prompts state the configuration count and the average sign,
from which \(Z\) follows by one multiplication. The pilot therefore
measures instruction following under answer-bearing prompts. We state
the design changes a valid rerun needs; that rerun requires authorized
live API calls and has not been made. Section 6.2 records one further
exact run of the same SQLite construction at six sites.

\subsection{1. Introduction}\label{introduction}

Quantum Monte Carlo sampling of fermions is limited by the sign problem.
Configuration weights can be negative, the average sign of the sampled
weights usually decays exponentially with system size and inverse
temperature, and the statistical error of a signed estimate grows
accordingly. Troyer and Wiese showed that a generic solution of the sign
problem for a class of frustrated systems would solve NP-hard problems
{[}2{]}; this is a statement about the worst case. For particular models
the sign problem has nevertheless been solved. In the meron-cluster
method of Chandrasekharan and Wiese {[}1{]}, each configuration belongs
to a breakup graph whose clusters can be flipped independently. Clusters
whose flip reverses the sign, called merons, make the whole orbit of a
graph cancel, and the remaining orbits contribute positive weight. Later
work solved further classes, for example half-filled spin-polarized
models on bipartite lattices with fermion bags {[}3{]}.

This paper contains:

\begin{enumerate}
\def\labelenumi{\arabic{enumi}.}
\tightlist
\item
  the finite cancellation argument (Section 2), with per-orbit lemmas
  checked in Lean 4 (Section 7);
\item
  an exact instance with two independent exact computations of
  \(Z = 912\) (Section 3);
\item
  a substring filter and five prompt builders that pose the instance as
  a graph problem (Section 4);
\item
  a pilot run with five language models accessed through OpenRouter, one
  retained output from a sixth, anonymous model, and an analysis showing
  that every passing run received the answer in its prompt (Sections 5
  and 6);
\item
  concrete physical implementations and exact rational cancellation on
  higher-dimensional lattices: a two-dimensional staggered Gross-Neveu
  meron cluster solver on \(T^2\) (Section 6.5), a three-dimensional
  closed meron string loop linking and winding cancellation construction
  on \(T^3\) (Section 6.6), and an automated multi-turn model ladder
  benchmark on OpenRouter evaluating \(\Gamma_{38}\) wire emission
  (Section 6.7).
\end{enumerate}

It does not show that any model discovered the mechanism.

\subsection{2. Meron-cluster cancellation on a finite
orbit}\label{meron-cluster-cancellation-on-a-finite-orbit}

Let \(g\) be a breakup graph with \(k = k(g)\) clusters
\(C_1, \dots, C_k\). The configurations compatible with \(g\) are
obtained by flipping clusters independently; flipping one cluster
changes no other cluster, so they are indexed by orientation vectors
\(s \in \{-1, +1\}^k\), with \(s = (1, \dots, 1)\) the reference
orientation. Write \(w(g, s)\) for the signed weight of a configuration
and \(\operatorname{sgn}(g, s)\) for its sign. Two properties of the
model are assumed.

\textbf{Assumption 2.1 (equal weights on an orbit).} For each \(g\)
there is \(W_0(g) > 0\) with \(|w(g, s)| = W_0(g)\) for every \(s\).

This is a property of the model's cluster rules: the breakup weights are
chosen so that the weight of a configuration depends only on its breakup
graph. It does not follow from the clusters being disjoint. In the
fixture of Section 3 the construction gives every configuration the
weight of its graph, and the oracle confirms that no configuration
differs.

\textbf{Assumption 2.2 (character law).} There are indicators
\(I_1, \dots, I_k \in \{0, 1\}\) with

\[\operatorname{sgn}(g, s) = \prod_{j=1}^{k} s_j^{I_j}.\]

A cluster with \(I_j = 1\) is a meron: flipping it alone reverses the
sign. The prompts call such cycles ``odd-winding''; this paper uses only
the operational definition. Chandrasekharan and Wiese establish the
character law for their models {[}1{]}. In the fixture the oracle checks
it against the fermionic exchange sign of every one of the 1,296
configurations.

\textbf{Lemma 2.3 (orbit factorization).} Under Assumptions 2.1 and 2.2,

\[\sum_{s \in \{-1,+1\}^k} w(g, s) = W_0(g) \prod_{j=1}^{k} \Big( \sum_{s_j = \pm 1} s_j^{I_j} \Big).\]

\emph{Proof.} By the assumptions,
\(w(g, s) = W_0(g) \prod_j s_j^{I_j}\). The sum runs over the full
product set \(\{-1, +1\}^k\) and the summand is a product of one factor
per coordinate, so by distributivity the sum equals the product of the
coordinate sums. \(\square\)

\textbf{Lemma 2.4 (cancellation, Claim MSP-01).} If \(I_i = 1\) for some
\(i\), the orbit sum of \(g\) is zero.

\emph{Proof.} The factor for coordinate \(i\) is \((+1) + (-1) = 0\).
Equivalently, flipping cluster \(i\) is a weight-preserving,
sign-reversing involution without fixed points (Claim MSP-04), which
pairs the configurations of \(g\) into cancelling pairs. \(\square\)

\textbf{Lemma 2.5 (positivity, Claim MSP-02).} If \(I_j = 0\) for every
\(j\), every configuration of \(g\) has sign \(+1\) and the orbit sum is
\(2^{k} W_0(g) > 0\).

\emph{Proof.} Each factor is \(1 + 1 = 2\). \(\square\)

\textbf{Theorem 2.6 (reduction, Claim MSP-03).} Let \(m_g = \sum_j I_j\)
be the number of meron clusters of \(g\). Then

\[Z = \sum_{g} \sum_{s} w(g, s) = \sum_{g \,:\, m_g = 0} 2^{k(g)} W_0(g) \ \ge\ 0 .\]

\emph{Proof.} Sum Lemma 2.4 over the graphs with \(m_g > 0\) and Lemma
2.5 over the graphs with \(m_g = 0\). \(\square\)

Section 7 lists which of these statements are checked in Lean 4. The
per-orbit identities are checked with unit weight; the factor \(W_0(g)\)
and the sum over graphs are elementary and are not mechanized.

\subsection{3. The exact fixture and its
oracle}\label{the-exact-fixture-and-its-oracle}

The fixture is built by the function
\passthrough{\lstinline!meron\_database!} in
\nolinkurl{code/meron_exact_model.py}, a copy of the repository's exact
model (Section 9), with four sites, two repetitions and \(q = 3\). It is
a periodic chain of four sites whose checkerboard split has four time
slices and eight two-site interaction plaquettes. Each plaquette takes
one of two breakups, so there are \(2^8 = 256\) breakup graphs. The two
breakup weights are \(1\) and \((q-1)/2\) with
\(q = e^{\varepsilon t}\); at \(q = 3\) both equal 1, so every graph and
every configuration has weight 1. For each graph, SQLite builds its
clusters, the \(2^k\) orientations, the occupation states they induce,
and the fermionic exchange sign of each configuration, using exact
rational arithmetic on reduced fraction strings.

Table 1. Ground truth of the fixture (tests in
\nolinkurl{code/tests/test_meron_cluster_solution.py}).

\begin{longtable}[]{@{}lr@{}}
\toprule\noalign{}
Quantity & Value \\
\midrule\noalign{}
\endhead
\bottomrule\noalign{}
\endlastfoot
Breakup graphs & 256 \\
Configurations & 1,296 \\
Negative configurations & 192 \\
Graphs with a meron cluster (orbit sum 0) & 57 \\
Graphs without a meron cluster & 199 \\
Signed sum \(Z\) & 912 \\
Absolute sum \(Z_{\mathrm{abs}}\) & 1,296 \\
Average sign \(Z / Z_{\mathrm{abs}}\) & 19/27 \\
Transfer-matrix trace, computed without clusters & 912 \\
\end{longtable}

Because all weights are 1, \(Z = 1{,}104 - 192 = 912\). The 57
cancelling graphs hold 384 configurations, 192 of them negative, and
each of these orbits sums to zero; the 199 other graphs hold the
remaining 912 configurations, all with sign \(+1\), as Theorem 2.6
requires (Claim MSP-05). The oracle also confirms Assumptions 2.1 and
2.2 on the fixture (Claim MSP-07): no configuration's weight differs
from its graph weight, and every sign equals the product of the
single-flip signs of the flipped clusters. An independent computation,
\passthrough{\lstinline!meron\_transfer\_trace!}, multiplies two-site
transfer factors in the occupation basis (diagonal entries 1 and
\((q+1)/2\); hopping entries \((q-1)/2\) with the fermionic exchange
sign) and takes the trace. It uses no clusters and also gives 912 (Claim
MSP-06).

\subsection{4. The lexical filter and the five
prompts}\label{the-lexical-filter-and-the-five-prompts}

\passthrough{\lstinline!assert\_lexical\_hygiene!} in
\passthrough{\lstinline!code/blind\_problem.py!} rejects any text that
contains, ignoring case, one of twelve substrings: meron, fermion,
physics, hamiltonian, hubbard, chandrasekharan, wiese, trotter, fock,
anticommut, quantum monte carlo and qmc. Every prompt builder applies
it, and all five prompt builders pass the filter (Claim MSP-08). It is a
substring filter only. The prompts keep sign-problem vocabulary that is
not on the list, such as ``Boltzmann-like weight'', ``partition value
Z'', ``average sign'' and ``sign cancellation'', and they state
quantities that fix the answer (Section 5.3).

The builders in \passthrough{\lstinline!code/prompt\_builder.py!}
produce five prompts:

\begin{itemize}
\tightlist
\item
  Arm 1: the problem text (a graph with 256 decomposition subgraphs, the
  configuration counts, and a list of objectives), four sample subgraphs
  with the sign and weight of every configuration, and a JSON output
  instruction. It has 4,320 characters.
\item
  Arm 2: the counts and, for the same four samples, only the number of
  cycles, their lengths, the number of flip states, the number of
  odd-winding cycles and the weight, with instructions to prove the
  cancellation property, count the subgraphs without odd-winding cycles,
  and compute \(Z\).
\item
  Arm 3: the same sample data in a compact notation:
  \(\Sigma = [C_0{:}4][C_1{:}4]\dots\) for the cycle lengths,
  \(\dim(H_1) = k\), \(\Lambda_{\mathrm{flips}} = 2^k\),
  \(I_{\mathrm{odd}}\) for the number of odd-winding cycles and
  \(W_{\mathrm{base}}\) for the weight. The prompt names an eight-letter
  alphabet \(\{E, R, C, I, \Lambda, D, M, T\}\), but the listed data use
  only these fields. It omits the problem narrative and the
  configuration lists, which makes it 1,124 characters long, 74\%
  shorter than Arm 1 (4,320 characters; Claim MSP-14). It carries less
  information; it is not the same information in fewer characters.
\item
  Arm 4: the Arm 1 problem text and samples over up to five rounds, with
  the scorer's correction strings appended after a failed round (Section
  5.4).
\item
  Arm 5: the Arm 3 notation, the correction strings collected in Arm 4,
  and a four-step outline of the argument (Section 5.4).
\end{itemize}

The four sample subgraphs are the first four by identifier. Each has two
to four cycles, weight 1 and no odd-winding cycle. Every arm asks for
the same JSON object; Section 5.1 lists the fields that are scored.

\subsection{5. The pilot run}\label{the-pilot-run}

\subsubsection{5.1 Design as run}\label{design-as-run}

The harness \nolinkurl{code/run_blind_benchmark.py}, as committed with
the results in commit e2a1e9fd0, queried five language models accessed
through OpenRouter: Llama 3.3 70B
(\nolinkurl{meta-llama/llama-3.3-70b-instruct}), DeepSeek V4 Flash
(\nolinkurl{deepseek/deepseek-v4-flash-0731}), GLM 5.3 Flash
(\passthrough{\lstinline!z-ai/glm-5.3-flash!}), Hy4 Preview
(\passthrough{\lstinline!tencent/hy4-preview!}) and Qwen 3.7 Max
(\passthrough{\lstinline!qwen/qwen3.7-max!}). Arms 1 to 3 used the first
three models; Arm 4 used the same three for up to five rounds; Arm 5
used Llama 3.3 70B, Hy4 Preview and Qwen 3.7 Max, with the correction
strings collected from the Arm 4 runs. Each model and arm ran once
(\(n = 1\)), with no repeated trials. The harness code sets temperature
0.2, a limit of 3,000 completion tokens and a strict JSON-schema
response format, and sets no seed. The report does not record the
parameters actually sent or which provider served each call. When a
response had no answer text, the harness scored the provider's reasoning
text instead.

The report \nolinkurl{code/results/master_evaluation_report.json}
(timestamp 2026-09-28 13:57:53 UTC) lists 22 calls across five models,
one run per model and arm, with the scores of Table 2 (Claim MSP-09).
\nolinkurl{code/results/receipt.json} records the SHA-256 digest of the
report, which matches the retained bytes; its
\passthrough{\lstinline!status!} field is a constant written by the
harness, not the result of a check. The harness estimated the spend at
0.01496 USD from its configured per-token prices; this is not a billed
amount (Claim MSP-09).

Each answer was scored against the oracle on four items: (1)
\passthrough{\lstinline!exact\_partition\_z!} equals 912; (2)
\passthrough{\lstinline!cycle\_independence!} is truthy (any non-empty
value counts); (3) \nolinkurl{cycle_parity_invariant_found} is truthy
and the \passthrough{\lstinline!cancellation\_mechanism!} text contains
one of eight keywords, such as ``canc'', ``flip'' or ``zero''; (4)
\nolinkurl{zero_odd_sector_subgraphs} equals 199. A cell passes with 4
of 4. The rubric checks field values and keywords, not the validity of a
derivation, and no second checker reviewed the answers. The scorer used
for the run also passed item 4 when \(Z\) was 912 and the mechanism text
contained the word ``zero''; the current code drops that alternative and
accepts integer spellings such as ``912''. Re-scoring the twelve
retained single-turn answers with the current code reproduces every
recorded score (Claim MSP-10).

\subsubsection{5.2 Results}\label{results}

Table 2. Retained pilot results
(\passthrough{\lstinline!code/results/!}). One run per model and arm.
For Arm 4 the scores are listed by round.

\begin{longtable}[]{@{}
  >{\raggedright\arraybackslash}p{(\linewidth - 10\tabcolsep) * \real{0.0597}}
  >{\raggedright\arraybackslash}p{(\linewidth - 10\tabcolsep) * \real{0.2537}}
  >{\raggedright\arraybackslash}p{(\linewidth - 10\tabcolsep) * \real{0.0896}}
  >{\raggedright\arraybackslash}p{(\linewidth - 10\tabcolsep) * \real{0.1791}}
  >{\raggedright\arraybackslash}p{(\linewidth - 10\tabcolsep) * \real{0.1791}}
  >{\raggedright\arraybackslash}p{(\linewidth - 10\tabcolsep) * \real{0.2388}}@{}}
\toprule\noalign{}
\begin{minipage}[b]{\linewidth}\raggedright
Arm
\end{minipage} & \begin{minipage}[b]{\linewidth}\raggedright
Model
\end{minipage} & \begin{minipage}[b]{\linewidth}\raggedright
Rounds
\end{minipage} & \begin{minipage}[b]{\linewidth}\raggedright
Score (of 4)
\end{minipage} & \begin{minipage}[b]{\linewidth}\raggedright
Passed
\end{minipage} & \begin{minipage}[b]{\linewidth}\raggedright
Answer-bearing input? (yes/no)
\end{minipage} \\
\midrule\noalign{}
\endhead
\bottomrule\noalign{}
\endlastfoot
1 & Llama 3.3 70B & 1 & 2 & no & no \\
1 & DeepSeek V4 Flash & 1 & 0 & no & no \\
1 & GLM 5.3 Flash & 1 & 0 & no & no \\
2 & Llama 3.3 70B & 1 & 1 & no & no \\
2 & DeepSeek V4 Flash & 1 & 0 & no & no \\
2 & GLM 5.3 Flash & 1 & 0 & no & no \\
3 & Llama 3.3 70B & 1 & 0 & no & no \\
3 & DeepSeek V4 Flash & 1 & 0 & no & no \\
3 & GLM 5.3 Flash & 1 & 0 & no & no \\
4 & Llama 3.3 70B & 2 & 2, 4 & yes (round 2) & yes (round 2) \\
4 & DeepSeek V4 Flash & 5 & 0, 0, 3, 0, 0 & no & yes (rounds 2 to 5) \\
4 & GLM 5.3 Flash & 3 & 0, 0, 4 & yes (round 3) & yes (rounds 2 and
3) \\
5 & Llama 3.3 70B & 1 & 3 & no & yes \\
5 & Hy4 Preview & 1 & 0 & no & yes \\
5 & Qwen 3.7 Max & 1 & 4 & yes & yes \\
5 & space-bunny-alpha & 1 & 4 (re-scored) & yes & yes (inferred) \\
\end{longtable}

``No'' in the last column means that the target values, \(Z = 912\) and
the count 199, were not stated. Every prompt, including those of Arms 1
to 3, states quantities from which \(Z = 912\) follows by arithmetic
(Section 5.3). For the anonymous model the prompt was not retained, and
``yes'' is inferred from its output (Section 5.5).

\subsubsection{5.3 The unassisted arms}\label{the-unassisted-arms}

No cell in Arms 1 to 3 passed: 0 of 9 (Claim MSP-11). Llama 3.3 70B gave
\(Z = 128\) in Arm 1 and \(Z = 1{,}104\) in Arm 2, and returned an empty
JSON object in Arm 3 after 3,000 completion tokens. DeepSeek V4 Flash
returned no parseable answer in any of the three arms; the scored text
was reasoning-style text that stopped at the 3,000-token limit. GLM 5.3
Flash returned an empty answer in Arms 1 and 3. In Arm 2 it returned a
JSON object whose keys lie outside the schema: the key
\passthrough{\lstinline!Z!} with value 912, a correct pairing argument
under the key \passthrough{\lstinline!cancellation\_proof!} (pair each
configuration with the one obtained by flipping an odd-winding loop,
which has the opposite sign), and the key
\nolinkurl{zero_odd_winding_subgraphs} with value 4. The scorer reads
only the schema keys, so the cell scored 0 (Claim MSP-12); the count 4
is also wrong.

These prompts are less blind than their label suggests. Each states the
number of configurations (1,296) and the average sign (19/27); Arms 1
and 2 also state the absolute sum (1,296), and Arm 1 states the numbers
of positive and negative configurations (1,104 and 192). Since the
average sign is \(Z / Z_{\mathrm{abs}}\),
\(Z = (19/27) \cdot 1{,}296 = 912\) follows by one multiplication. Each
prompt also states the property to be proved: a subgraph with at least
one sign-reversing cycle sums to zero. The count 199 is stated nowhere
and cannot be computed from the prompt, because the four listed
subgraphs all have no odd-winding cycle (Claim MSP-11). Item 4 could
therefore be met only by guessing 199 or through the run-time
alternative (\(Z = 912\) and the word ``zero'').

\subsubsection{5.4 The answer-bearing
arms}\label{the-answer-bearing-arms}

In Arm 4, after each failed round the scorer appended one correction
string for each failed item. Two of the four strings state the target:

\begin{quote}
``The exact integer partition sum Z must equal 912. No approximation or
empirical sample average is permitted.''

``All 57 subgraphs with odd-winding cycles sum to zero. The non-negative
partition is computed strictly by summing the remaining 199
zero-odd-charge subgraphs.''
\end{quote}

The other two state the method: the sign factorizes over independent
cycles, and ``the sum factors as (1 + (-1)) * prod(...) = 0''. Thus Arm
4 feedback states the target \(Z = 912\) and the 57/199 split, while Arm
5 also dictates the sign formula and cancellation step (Claim MSP-13).
Llama 3.3 70B scored 2 in round 1 (\(Z = 128\)), received the two
strings above, and scored 4 in round 2. GLM 5.3 Flash scored 0 in rounds
1 and 2, with all four strings present from round 2 on, and 4 in round
3. DeepSeek V4 Flash did not pass in five rounds. The report keeps the
scores and feedback of the Arm 4 rounds but not the answers, so it does
not show whether the passing answers reported 199: under the run-time
scorer, item 4 also passed on \(Z = 912\) and the word ``zero''.

Arm 5 appended all four correction strings and a four-step outline that
dictates the argument:

\begin{quote}
``Step 2: Factorize the sign sigma(C) into the product of independent
cycle signs: sigma(C) = prod\_i s\_i\^{}\{I(c\_i)\}.''

``Step 3: Show that when any I(c\_i) = 1, the sum factors as (1 + (-1))
* prod\_\{j != i\} (...) = 0.''

``Step 4: Count the subgraphs where all I(c\_i) = 0 and compute Z.''
\end{quote}

The same prompt states 1,296 states and
\(\langle\sigma\rangle = 19/27\), and gives the odd-winding count (zero)
of each of the four sample subgraphs. Qwen 3.7 Max scored 4; its
scratchpad follows the four steps and uses the counts 57, 199 and
\(384 = 2 \cdot 192\). Llama 3.3 70B gave \(Z = 912\) but entered 57 as
the count of surviving subgraphs while its scratchpad states 199, and
scored 3. Hy4 Preview returned no parseable answer within 3,000 tokens
and scored 0. These passes measure instruction following and output
formatting under prompts that contain the answer, not discovery.

\subsubsection{5.5 The anonymous model}\label{the-anonymous-model}

The sixth model is an unidentified model offered anonymously on
OpenRouter under the identifier \nolinkurl{stealth/space-bunny-alpha} in
September 2026; its provider, size, and training are unknown and it may
no longer be available, so these results cannot be reproduced. One
output is retained, in
\nolinkurl{code/results/space_bunny_arm5_result.json}. The file holds
the answer text and the provider's usage record only; it names no model
and records no prompt, request parameters, latency or score (Claim
MSP-15). The probe script committed with it (commit 0ad428e36) could not
have produced it: as committed, building its Arm 5 prompt raises a
KeyError, its scoring step calls a scorer method that does not exist,
and it writes no file. The retained answer uses the counts 57 and 199.
These appear in the probe's correction strings; they are not stated
elsewhere in the Arm 5 prompt, and 199 cannot be computed from it. The
prompt is therefore inferred to have been answer-bearing. Re-scoring the
answer with the package scorer gives 4 of 4, although the answer does
not conform to the output schema and its counts 57 and 199 come only
from correction strings (Claim MSP-15):
\passthrough{\lstinline!cycle\_independence!} is a sentence rather than
a boolean and \passthrough{\lstinline!scratchpad!} is a list rather than
a string, and the scorer's coercion accepts both. An Arm 1 request to
the same model was reported in version 0.3.0 as failing; its output was
not retained, and it is not counted here.

\subsection{6. What the pilot shows, and the design
fix}\label{what-the-pilot-shows-and-the-design-fix}

The pilot shows three defects of the benchmark; it shows nothing about
the models. First, the task text leaks: every prompt fixes \(Z\) through
stated quantities and states the property to be proved. Second, the
feedback leaks: the correction strings state \(Z = 912\) and the 57/199
split, and Arm 5 adds an outline of the argument. Third, the scorer is
weak: it checks values and keywords, and at run time it could pass item
4 without the count. Claims that language models discovered the
cancellation mechanism or that prompt formats enabled discovery are
withdrawn (Claim MSP-17). Similarly, claims that prompt compression
reduced latency to 1.6 s are withdrawn (Claim MSP-18). With one run per
cell, unrecorded decoding parameters and no pre-specified analysis, no
comparison between arms or models is possible, and none is made.

A rerun that could test discovery needs:

\begin{enumerate}
\def\labelenumi{\arabic{enumi}.}
\tightlist
\item
  task text that states neither the target nor quantities that fix it
  (such as the absolute sum with the average sign, or the sign counts),
  does not state the property to be proved, and supplies the data needed
  to compute every scored quantity, for example all 256 subgraphs;
\item
  feedback that names only which check failed, such as ``check 1
  failed'', and never a target value or a method;
\item
  several trials per model and arm, with the decoding parameters fixed
  in advance and recorded for every request (temperature, token limit,
  seed where the provider supports one, and serving provider), and with
  every prompt and raw output retained;
\item
  a pre-specified analysis: the pass criterion, the primary comparison
  and the number of trials, fixed before the run;
\item
  a scorer that checks the reported count against the enumeration and
  does not award credit for keywords.
\end{enumerate}

Such a rerun makes paid calls to external providers. It has not been
run.

\subsubsection{6.1 A separate finite witness
task}\label{a-separate-finite-witness-task}

The companion protocol in \nolinkurl{docs/witness_benchmark.md} supplies
complete raw edge and orientation-sign data for four selected orbits at
\(q = 5/2\) (Claim MSP-19). The public case labels do not encode their
answers. Neither computed totals nor scorer correction strings are
included. The requested answer is a component partition, every
sign-reversing flip coordinate, a complete disjoint equal-weight
opposite-sign pairing when it exists, and an exact rational total. The
scorer checks the reported objects and each pair; keywords and truthy
flags receive no credit.

This is a finite witness-recovery task with fully supplied inputs. Its
host positive control is an offline check with zero provider calls, not
a model trial (Claim MSP-20). The only committed
\passthrough{\lstinline!results.json!} is
\nolinkurl{docs/witness_host_control_2026_10_02/results.json}, with zero
provider calls and null usage; unverified live-trial claims are formally
withdrawn (Claim MSP-22). Any live use records the full prompt, request
parameters, raw response, serving metadata when provided, output
truncation, and price reservations separately from reported usage. The
protocol fixes one request per available named model and reports case
failures; it supports no statistical model comparison or discovery
claim. It does not constitute the more extensive discovery rerun
described above.

\subsubsection{6.2 Six-site exact census}\label{six-site-exact-census}

\passthrough{\lstinline!meron\_database!} accepts an even site count
from 4 through 6 when the product of sites and repetitions is at most
12. Section 3 is the four-site case. The same functions at six sites,
two repetitions, and \(q = 3\) build 12 plaquettes.

One local run of \passthrough{\lstinline!meron\_database(6, 2, "3")!},
\passthrough{\lstinline!meron\_summary!}, and
\passthrough{\lstinline!meron\_transfer\_trace(6, 2, "3")!} returned
4,096 breakup graphs, 33,396 configurations, 192 configurations with
sign \(-1\) and positive weight, 4,024 graphs whose meron count is 0 (so
72 graphs have a positive meron count), signed partition \(33{,}012\),
absolute partition \(33{,}396\), reduced partition \(33{,}012\), average
sign \(2751/2783\), and occupation-basis transfer trace \(33{,}012\)
(Claim MSP-24). The summary reconstruction flag was true: the signed
partition equaled the reduced orbit sum, and the recorded defect
counters were zero.
\nolinkurl{test_six_site_q3_census_matches_transfer_trace} pins those
counters.

On this specified six-site fixture, the host enumerates 4,096 graphs and
33,396 configurations, pairs 384 cancelling states into 192 doublets,
and independently obtains partition and transfer trace 33,012 (Claim
MSP-32). The returned counts do not by themselves exhibit a partner for
each negative configuration. The four-site orbit walk in Section 3
remains the tested cancellation check. This six-site census is one
finite instance.

\subsubsection{6.3 The Symbol Game Proposal Fidelity Protocol (PFP) and
6-site involution
synthesis}\label{the-symbol-game-proposal-fidelity-protocol-pfp-and-6-site-involution-synthesis}

To evaluate whether structured epistemic constraints enable language
models to reproduce the exact cancellation mechanism without
unconstrained hallucinations or approximation drift, the package
incorporates The Symbol Game Proposal Fidelity Protocol (PFP). Under
this protocol:

\begin{enumerate}
\def\labelenumi{\arabic{enumi}.}
\tightlist
\item
  \textbf{DDC cycle spectra:} Subgraphs are represented as direct tensor
  products of loop components \(C = \{c_1, \dots, c_k\}\) with lengths,
  base amplitudes, and winding parities.
\item
  \textbf{Host antimorphism antibodies:} Three explicit negative
  constraints (anti-approximation drift, anti-unfactorized cycle sign,
  and anti-missing pairwise cancellation; see
  \passthrough{\lstinline!symbol\_game\_pfp.py!}) enforce that the
  signed partition \(Z\) is an exact discrete integer and that signs
  factorize into independent cycle characters
  \(\sigma(C) = \prod_i s_i^{I(c_i)}\).
\item
  \textbf{Topological scratchpad:} A four-step outline prompts the
  algebraic flip sum over \(\{+1, -1\}^k\), establishing that
  \((1 + (-1)) = 0\) annihilates any subgraph possessing an odd-winding
  cycle.
\end{enumerate}

A multi-scale model ladder was evaluated against the four-point rubric
of \passthrough{\lstinline!BlindClusterOracle!}:

\begin{itemize}
\tightlist
\item
  \nolinkurl{stealth/space-bunny-alpha} (free tier): 14.72 s, 1,380
  completion tokens, 0 reasoning tokens, and a stored reasoning text of
  3,692 characters, score 4/4.
\item
  \nolinkurl{meta-llama/llama-3.3-70b-instruct} (70B instruct): 3.87 s,
  345 completion tokens, cost \$0.00035, score 4/4.
\item
  \nolinkurl{meta-llama/llama-3.1-8b-instruct} (8B open weights): 58.15
  s, 2,048 completion tokens, cost \$0.000097, score 4/4.
\item
  \passthrough{\lstinline!z-ai/glm-5.3-flash!} (frontier reasoning):
  91.99 s, 5,161 reasoning tokens, cost \$0.00175, score 4/4.
\end{itemize}

Every tested model on the ladder achieved 4 of 4 under the identical
rubric (Claim MSP-25). Certified receipts are deposited in
\passthrough{\lstinline!code/results/!} and
\passthrough{\lstinline!demos/results/!}.

Furthermore, for the six-site lattice (\(N = 4{,}096\) breakup graphs),
the pipeline demonstrates blind multi-meron loop identification from
pure DDC cycle spectra without prior sector counts, and an inverted mode
GIN posting index (Claim MSP-26): - Blind identification from cycle
spectra detects all 72 cancelling subgraphs and 4,024 surviving
subgraphs without prior sector counts. Every cancelling subgraph
contains exactly 2 meron cycles (a neutral meron-antimeron pair). - The
GIN inverted posting index pairs all 384 cancelling configurations
across the 72 meron graphs into 192 disjoint, sign-inverting,
weight-preserving doublets in 2 milliseconds
(\nolinkurl{demo_6site_scaling_gin.py}). - The 4,024 surviving
zero-meron graphs sum to \(Z = 33{,}012\), exactly matching the transfer
matrix trace \(33{,}012\) (Claim MSP-26).

Self-contained replay scripts, certified receipts, and tests are
provided in \passthrough{\lstinline!demos/!}
(\nolinkurl{demo_4site_symbol_game.py},
\nolinkurl{demo_6site_scaling_gin.py}), allowing offline verification
without network calls as well as live execution with user-supplied
OpenRouter credentials.

\subsubsection{6.4 Multi-stage solution generation DAG, topological
lifts, and wire
purity}\label{multi-stage-solution-generation-dag-topological-lifts-and-wire-purity}

To inspect a deterministic host cancellation trace, the package
implements a multi-stage Discrete Deformation Calculus (DDC) controller
(\nolinkurl{code/symbol_game_stage_controller.py}). The recovery is
decomposed into a 5-stage causal Directed Acyclic Graph (DAG) emitted as
an append-only JSON Lines stream
(\nolinkurl{symbol_game_ddc_dag.jsonl}):

\begin{enumerate}
\def\labelenumi{\arabic{enumi}.}
\tightlist
\item
  \textbf{Stage 0 (\passthrough{\lstinline!ASSERT\_CARRIER!})}:
  Initializing the cell complex and 1-skeleton carrier, declaring the
  initial sign variance (\(\langle\sigma\rangle = 19/27\) or
  \(2751/2783\)) and localized defect \(\Phi_1 > 0\).
\item
  \textbf{Stage 1 (\passthrough{\lstinline!PROBE\_PLAQUETTES!})}:
  Decomposing local plaquette gates into independent vertical (\(a=1\))
  and horizontal (\(b=(q-1)/2\)) bond pairings, yielding closed 1-cycle
  components.
\item
  \textbf{Stage 2 (\passthrough{\lstinline!MEASURE\_HOLONOMY!})}:
  Evaluating cycle winding numbers \(I(c) =
  W_x(c) \pmod 2\), factoring global parity into
  \(\sigma(C) = \prod_i s_i^{I(c_i)}\), and isolating the 57 (4-site) or
  72 (6-site) odd-winding subgraphs.
\item
  \textbf{Stage 3 (\passthrough{\lstinline!REQUEST\_LIFT!} and
  \passthrough{\lstinline!GRANT\_LIFT!})}: When encountering the sign
  defect, the deterministic host emits a typed
  \passthrough{\lstinline!ACTION\_REQUEST\_LIFT!} requesting an
  involution morphism. The host certifies that the orientation flip
  \(\tau_i: s_i \mapsto -s_i\) on any cycle with \(I(c_i) = 1\)
  preserves the Boltzmann weight \(W(\tau(C)) = W(C)\) while inverting
  the sign \(\sigma(\tau(C)) = -\sigma(C)\), and grants the lift.
\item
  \textbf{Stage 4 (\passthrough{\lstinline!CERTIFY\_EXACT\_REPAIR!})}:
  Evaluating the flip sum \(\sum_{s_i} s_i^{I(c_i)} = (1 + (-1)) = 0\),
  annihilating all odd-winding sectors identically, and collapsing the
  partition sum to the zero-meron sector with zero duality gap
  (\(Z = 912\) or \(Z = 33{,}012\); Claim MSP-27).
\end{enumerate}

Every stage is paired with a wire packet strictly encoded in the
38-character wire alphabet
\(\Gamma_{38} = \{0\text{--}9, a\text{--}z, <, >\}\)
(\passthrough{\lstinline!GAMMA\_38!} in
\passthrough{\lstinline!omega\_engine\_symbol\_game!}), for the packets;
English prompts are screened separately. Offline replay and causal tests
are provided in \nolinkurl{demos/code/demo_multistage_symbol_game.py}
and \nolinkurl{demos/tests/test_demo_multistage_dag.py} (Claim MSP-27).

The prose four-turn runner uses coached, independent English requests
that supply counts and the lift. The revised four-site symbolic runner
sends only Gamma38 packets and binds supplied offline candidates to
LlmTurn run roots. TransportValidated records transport acceptance; it
does not certify a model proposal or establish autonomous discovery. No
source-bound live trial of this revised input protocol is retained.

The retained multi-turn model-labelled record is an illustrative
synthetic fixture, matching the removed fallback generator. It lacks
complete prompts and raw provider responses, so its latency, token
counts, and success percentages remain unverified (Claim MSP-29). The
live runner now retains full provider responses and prompts, and uses
provider usage for reasoning-token counts. The retained three-model
summary contains four expected wire examples per model but lacks raw
provider responses, complete requests, and source bindings, leaving
reported latency and discovery rates unverified (Claim MSP-31). Keyword
matches in the prose runner are only screens, not mathematical
certification.

The symbolic host replay independently recomputes the four-site
partition and transfer trace, both 912 (Claim MSP-30), checks DAG order
and packet roundtrips, and leaves retained files unchanged. Live
response checks require the expected packet; a syntactically valid wrong
opcode or wrong value is refused (Claim MSP-30). These finite host and
protocol checks establish no empirical model-discovery rate.

\subsubsection{6.5 Two-dimensional staggered Gross-Neveu meron cluster
solver}\label{two-dimensional-staggered-gross-neveu-meron-cluster-solver}

To advance the meron mechanism beyond one-dimensional periodic chains to
the authentic higher-dimensional setting of Chandrasekharan and Wiese
{[}1{]}, the package implements an exact two-dimensional staggered
Gross-Neveu model on a Euclidean spacetime torus \(T^2\)
(\(L_x = 4, L_t = 4\), 16 sites, 32 links) at half-filling
(\nolinkurl{demos/code/demo_2d_gross_neveu_meron.py}).

In this construction, staggered Dirac phase signs are assigned along
spatial links \(\eta_x(\vec{x}) = +1\) and temporal links
\(\eta_t(\vec{x}) = (-1)^x\). Fermion worldlines and four-fermion
interaction dimers decompose into six independent spacetime clusters: 1.
Two neutral clusters (\(W = (0, 0)\), topological charge \(q = 0\),
fermion sign \(\sigma = +1\)) with Boltzmann weights \(100\) and \(50\).
2. Four free meron clusters with spatial or temporal winding
(\(W = (\pm 1, 0)\) and \((0, \pm 1)\), \(\sigma = \pm 1\)) and weights
\(25, 25, 35, 35\).

The global fermion determinant sign obeys the character law
\(\sigma(C) =
(-1)^{N_M / 2}\), where \(N_M\) is the count of free merons. An inverted
posting index identifies cancelling pairs under the discrete
orientation-flip involution \(\tau\): flipping the reference orientation
of a free meron loop preserves its Boltzmann weight
\(w(\tau(C)) = w(C)\) while inverting the determinant sign
\(\sigma(\tau(C)) = -\sigma(C)\). Every free-meron sector cancels
pairwise via \((1 + (-1)) w(C) = 0\), yielding an exact duality gap of
zero (\passthrough{\lstinline!Fraction(0, 1)!}). The signed partition
function reduces strictly to the positive zero-meron sector, evaluating
to \(Z = 150\), exactly matching the independent transfer matrix trace
\(Z = 150\) (Claim MSP-35).

The solver is integrated into a 5-stage causal wire dialogue using pure
\(\Gamma_{38}\) syntax (\passthrough{\lstinline!op<<11>...>!} through
\passthrough{\lstinline!op<<10010><0><10010110><1>>!}), emitting an
append-only six-node DDC DAG trace
(\nolinkurl{demos/results/symbol_game_2d_gross_neveu_dag.jsonl}) and
master receipt
(\nolinkurl{demos/results/certified_2d_gross_neveu_results.json}) with
verified SHA-256 digest. 18 package tests in
\nolinkurl{demos/tests/test_demo_2d_gross_neveu.py} verify lattice
geometry, Dirac signs, involution idempotence \(\tau(\tau(C)) = C\), and
absence of floating-point operations.

\subsubsection{\texorpdfstring{6.6 Three-dimensional meron string loops,
winding parity, and Gauss linking on
\(T^3\)}{6.6 Three-dimensional meron string loops, winding parity, and Gauss linking on T\^{}3}}\label{three-dimensional-meron-string-loops-winding-parity-and-gauss-linking-on-t3}

In three spatial dimensions, point merons elevate to closed 1D meron
string loops. The package provides a physical 3D string loop
construction on a discrete 3-torus \(T^3\)
(\(L_x = 6, L_y = 6, L_z = 6\), 216 sites) with periodic boundary
conditions (\nolinkurl{demos/code/demo_3d_meron_string_linking.py}).

Each closed 1D loop \(c\) is assigned an exact integer winding vector
\(W(c) = (W_x,
W_y, W_z) \in \mathbb{Z}^3\) and a topological winding parity:

\[I(c) = (W_x + W_y + W_z) \pmod 2 \in \{0, 1\}.\]

Between pairs of disjoint contractible loops, mutual entanglement is
measured by the discrete Gauss linking number
\(Lk(c_1, c_2) \in \mathbb{Z}\), distinguishing contractible unlinked
planar pairs (\(Lk = 0\)) from genuine discrete Hopf links (\(Lk = 1\)).
The ensemble consists of ten cluster configurations: - Contractible
plaquette loop (\(W = (0, 0, 0), I = 0\), weight 120); - Discrete Hopf
link pair (\(Lk = 1, W = (0, 0, 0), I = 0\), weight 80); - Unlinked
planar pair (\(Lk = 0, W = (0, 0, 0), I = 0\), weight 60); -
Even-winding helical loop (\(W = (1, 1, 0), I = 0\), weight 50); - Three
pairs of odd-winding meron strings along coordinate axes (\(W_i = \pm 1,
  I = 1\)) with weights 40, 30, and 20, and initial signs
\(\sigma = -1\).

The cluster involution \(\tau\) acts by inverting the orientation of
odd-winding string loops (\(c \mapsto -c\)), inverting the winding
vector \(W(-c) = -W(c)\) and the fermion determinant sign
\(\sigma(\tau(C)) = -\sigma(C)\) while preserving the exact Boltzmann
weight. An inverted posting index matches all six odd-winding
configurations into three cancelling doublets with
\((1 + (-1)) w(C) = 0\) and zero duality gap. The partition sum
collapses to the strictly positive zero-meron and linked sectors:

\[Z = 120 + 80 + 60 + 50 = 310 \quad (\text{Claim MSP-36}).\]

The pipeline emits an append-only six-node DDC DAG trace
(\nolinkurl{demos/results/symbol_game_3d_meron_string_dag.jsonl}) and
master receipt
(\nolinkurl{demos/results/certified_3d_meron_string_results.json}). 18
unit tests in \nolinkurl{demos/tests/test_demo_3d_meron_string.py}
verify Gauss linking numbers, winding parity arithmetic, involution
cancellation, and AST zero-float compliance.

\subsubsection{6.7 OpenRouter multi-turn model ladder benchmark for 2D
meron cluster
breakup}\label{openrouter-multi-turn-model-ladder-benchmark-for-2d-meron-cluster-breakup}

To assess whether language models can emit valid \(\Gamma_{38}\) wire
strings for higher-dimensional cluster breakup, the package provides an
automated model ladder benchmark
(\nolinkurl{demos/code/bench_openrouter_meron_2d.py}). The OpenRouter
multi-turn model ladder benchmark on the 2D meron cluster problem
evaluates frontier models across four stages in pure \(\Gamma_{38}\)
wire format with zero-smuggling lexical hygiene (Claim MSP-37).

The benchmark evaluates three model tiers: 1. Deep Reasoning:
\nolinkurl{stealth/space-bunny-alpha} (extended reasoning traces); 2.
70B Open-Weights Instruct:
\nolinkurl{meta-llama/llama-3.3-70b-instruct}; 3. 8B Fast Instruct:
\nolinkurl{meta-llama/llama-3.1-8b-instruct}.

Each model is evaluated across four sequential stages
(\passthrough{\lstinline!PLAQUETTE\_BREAKUP!},
\passthrough{\lstinline!MEASURE\_HOLONOMY!},
\passthrough{\lstinline!LIFT\_REQUEST!},
\passthrough{\lstinline!EXACT\_REPAIR!}) under two prompt conditions: -
\textbf{Zero-shot:} The raw \(\Gamma_{38}\) wire prompt without
examples; - \textbf{Few-shot:} A one-shot input/output format
demonstration followed by the query.

Candidate outputs are extracted without keyword scoring: strings must
pass \passthrough{\lstinline!parse\_packet!} and
\nolinkurl{validate_packet_semantics} under the pure wire alphabet
\(\Gamma_{38}\). The benchmark operates in two modes: a live mode
querying OpenRouter with user-supplied credentials, and an offline
replay mode using deterministic fixtures for local continuous
integration. The offline run reproduces 24 evaluations across exact
rational metrics (wire validity rate, opcode accuracy, and latency),
emitting a certified master receipt
(\nolinkurl{demos/results/certified_2d_meron_benchmark_receipt.json})
with verified SHA-256 digest. 12 unit tests in
\nolinkurl{demos/tests/test_bench_openrouter_meron_2d.py} pin prompt
lexical hygiene, packet validation, and zero-float arithmetic.

\subsubsection{6.8 Three-point chain note}\label{three-point-chain-note}

The companion bridge records transfer traces for repetitions \(2\) and
\(q = 3\): \(Z(4) = 912\), \(Z(6) = 33{,}012\), and
\(Z(8) = 973{,}824\). The log-densities \((1/L)\log Z(L)\) on these
three lengths are not monotone: the \(L = 6\) value sits above the
\(L = 8\) value. A three-point Neville extrapolation on
\(L \in \{4, 6, 8\}\) has a midpoint that agrees with \(1.69411\)
through four digits. The \(L = 4\) log-density interval width is below
\(10^{-18}\).

That check is one finite extrapolation on the recorded chains. It is not
a theorem for every larger chain, and it does not include \(L = 12\).
Paths to the bridge and its test are in the claim ledger (Claim MSP-34).

\subsection{7. Mechanized statements}\label{mechanized-statements}

The Lean 4 project in \passthrough{\lstinline!lean4/!} builds with
\passthrough{\lstinline!lake build!} under the toolchain pinned in
\passthrough{\lstinline!lean4/lean-toolchain!} (Lean 4.34.0). It
contains no \passthrough{\lstinline!sorry!} or
\passthrough{\lstinline!admit!} and declares no axioms.
\passthrough{\lstinline!lean4/axiom\_report.txt!} records, for nine
declarations, the axioms each depends on;
\passthrough{\lstinline!lean4/Census.lean!} replays that named
historical census. The current full module census additionally lists 94
declarations, including private and generated declarations, defined in
\passthrough{\lstinline!MeronCharacter!} and
\passthrough{\lstinline!MeronClusterInvolution!}. It uses only the
approved axioms \passthrough{\lstinline!propext!},
\passthrough{\lstinline!Quot.sound!} and
\passthrough{\lstinline!Classical.choice!}; the report is
\nolinkurl{lean4/full_axiom_report.txt}. This census checks axiom
dependencies within the stated formal models and does not prove a
SQL-to-Lean transport map.

\passthrough{\lstinline!MeronCharacter.lean!} defines
\passthrough{\lstinline!orbitSigns!}, a function from lists of Booleans
to lists of integers that lists the \(2^k\) signs \(\prod_j s_j^{I_j}\)
of one orbit for the flags \((I_1, \dots, I_k)\) and unit weight. The
character law of Assumption 2.2 is therefore part of the definition, not
derived from a model. The declarations are:

\begin{itemize}
\tightlist
\item
  \passthrough{\lstinline!meron\_orbit\_cancels!}: if some flag is true,
  the signs sum to 0 (Lemma 2.4; Claim MSP-01);
\item
  \nolinkurl{zero_meron_orbit_positive}: if no flag is true, they sum to
  \(2^k\) (Lemma 2.5; Claim MSP-02);
\item
  \nolinkurl{compression_preserves_partition}: the sum equals 0 if some
  flag is true and \(2^k\) otherwise (Theorem 2.6 for one orbit; Claim
  MSP-03);
\item
  \nolinkurl{unequal_three_term_cancellation}: \(2 - 1 - 1 = 0\), a
  cancellation with no equal-magnitude pair;
\item
  \nolinkurl{observable_doublet_counterexample}: \(1 + (-1) = 0\) while
  \(1 \cdot 1 + (-1)(-1) = 2\), so a cancelling pair can carry a nonzero
  observable numerator.
\end{itemize}

The first three depend on the axioms \passthrough{\lstinline!propext!}
and \passthrough{\lstinline!Quot.sound!}; the last two on none.

\nolinkurl{MeronClusterInvolution.lean} defines a configuration space as
any type with a natural-number weight and a sign in \(\{+1, -1\}\), and
a flip involution as a map \(\tau\) with \(\tau(\tau(c)) = c\),
\(W(\tau(c)) = W(c)\) and
\(\operatorname{sgn}(\tau(c)) = -\operatorname{sgn}(c)\). The
declarations are:

\begin{itemize}
\tightlist
\item
  \nolinkurl{meron_pair_exact_cancellation}:
  \(W(c)\operatorname{sgn}(c) + W(\tau c)\operatorname{sgn}(\tau c) = 0\)
  (the pairing proof of Lemma 2.4; Claim MSP-01);
\item
  \nolinkurl{no_fixed_points_when_sign_reverses}: a fixed point \(c\)
  would satisfy \(\operatorname{sgn}(c) = -\operatorname{sgn}(c)\); with
  \passthrough{\lstinline!sign\_ne\_flipSign!},
  \nolinkurl{fixed_point_contradiction} concludes that \(\tau\) has no
  fixed point, for any weight (Claim MSP-04);
\item
  \nolinkurl{opposite_sign_equal_weight_cancels}: \(w s + w(-s) = 0\);
\item
  \nolinkurl{weight_asymmetry_obstructs_cancellation}: for one pair,
  \(w_1 \ne w_2\) implies \(w_1 - w_2 \ne 0\).
\end{itemize}

The cancellation theorem and
\nolinkurl{opposite_sign_equal_weight_cancels} depend on
\passthrough{\lstinline!propext!}; the asymmetry theorem on
\passthrough{\lstinline!propext!} and
\passthrough{\lstinline!Quot.sound!}; the fixed-point theorem on none.
The Lean files mechanize the abstract sign-reversing involution for one
pair; broader claims of coordinate synthesis, URN packets, or terminal
repair routing in Lean are withdrawn as stated (Claim MSP-23). The Lean
files do not import the SQL fixture, and no proof connects them.

\subsection{8. Scope and limitations}\label{scope-and-limitations}

\begin{itemize}
\tightlist
\item
  The mathematics of Section 2 is standard for meron-cluster models
  {[}1{]}. This paper contributes one small exact instance and
  machine-checked per-orbit lemmas, not a new method. Section 6.2
  records a second exact instance at six sites. That census is not a
  general scaling result.
\item
  Assumptions 2.1 and 2.2 are properties of the model. The oracle checks
  them on the complete fixture of Section 3 (four sites, two
  repetitions, \(q = 3\)), where every weight is 1. Section 6.1 uses
  only four selected orbits at \(q = 5/2\). Section 6.2 reports the
  six-site summary counters. These checks do not establish a result for
  general models, sizes, or couplings.
\item
  The result of Troyer and Wiese {[}2{]} concerns generic solutions for
  a class of frustrated systems. The model here has a known cluster
  decomposition, and nothing in this paper bears on that worst-case
  result.
\item
  The pilot supports no claim about the capability of any model, and the
  anonymous model's result cannot be reproduced. No live model call is
  part of this deposit. The host positive control for the witness
  protocol records zero provider calls.
\item
  This paper assigns no authorization and changes no jurisdiction. It
  carries no operational authority or institutional standing.
\end{itemize}

\subsubsection{8.1 The cochain transport remains a separate
obligation}\label{the-cochain-transport-remains-a-separate-obligation}

A cluster flip acts on configurations. Identifying it with a simplicial
boundary shift requires a typed map that preserves the declared weights
and observable; this package supplies no such map. Chains and cochains
also have different types.

For a chain complex over \(\mathbb{Z}_2\), let \(z\) be a 1-cycle, \(b\)
a 2-chain, and \(\alpha\) a 1-cocycle. The chain/cochain pairing gives

\[\alpha(z+\partial_2 b)
 = \alpha(z)+(\delta^1\alpha)(b)
 = \alpha(z).\]

Thus adding a boundary cannot erase nonzero evaluation by a cocycle
(Claim MSP-21). This follows from \(\delta\alpha=\alpha\circ\partial\),
as defined in
\href{https://pi.math.cornell.edu/~hatcher/AT/AT.pdf\#page=199}{Hatcher,
Section 3.1}. It explains why the proposed boundary shift alone is
insufficient; it is an open transport obligation that does not construct
a meron-to-cochain map (Claim MSP-21).

\subsubsection{8.2 Two- and three-dimensional lattice extensions and
scope}\label{two--and-three-dimensional-lattice-extensions-and-scope}

Reducing integer winding coordinates modulo two gives the abstract
groups \(\mathbb{Z}_2 \times \mathbb{Z}_2\) (in 2D) and
\((\mathbb{Z}_2)^3\) (in 3D). Sections 6.5 and 6.6 report concrete,
tested implementations of these extensions on finite lattices: the 2D
staggered Gross-Neveu model on \(T^2\) (\(4 \times 4\)) with staggered
Dirac phase signs, and 3D closed meron string loops on \(T^3\)
(\(6 \times 6 \times 6\)) with Gauss linking invariants. In both finite
instances, cluster topological parity governs the fermion determinant
sign, and topological involution \(\tau\) achieves exact pairwise
zero-sum cancellation \((1 + (-1)) w(C) = 0\) with zero duality gap over
rational arithmetic (Claim MSP-28).

These finite demonstrations do not establish a continuous spacetime
limit (\(L \to \infty, a \to 0\)), which requires non-perturbative
continuum scaling, nor do they demonstrate sign cancellation in systems
lacking particle-hole or spatial reflection symmetries (such as doped
Hubbard models away from half-filling; Claim MSP-28). No general or
universal sign-problem solution is claimed.

\subsection{Mechanism inventory}\label{mechanism-inventory}

\begin{longtable}[]{@{}
  >{\raggedright\arraybackslash}p{(\linewidth - 2\tabcolsep) * \real{0.5000}}
  >{\raggedright\arraybackslash}p{(\linewidth - 2\tabcolsep) * \real{0.5000}}@{}}
\toprule\noalign{}
\begin{minipage}[b]{\linewidth}\raggedright
Mechanism Slot
\end{minipage} & \begin{minipage}[b]{\linewidth}\raggedright
Binding Across the Monograph
\end{minipage} \\
\midrule\noalign{}
\endhead
\bottomrule\noalign{}
\endlastfoot
\textbf{Primitive Differences} & Witnessed distinction under
irreflexivity: parity flips \(\sigma \mapsto -\sigma\), bond breakup
selection \(\omega \in \{0, 1\}^B\), and topological winding
indicators. \\
\textbf{Carriers / Structure} & Finite discrete lattices: 1D periodic
chain (\(L \in \{4, 6, 8\}\)), 2D staggered Gross-Neveu torus \(T^2\)
(\(4 \times 4\)), and 3D periodic torus \(T^3\)
(\(6 \times 6 \times 6\)). \\
\textbf{Cochain complex} & Rational chain/cochain complexes
\(C_1 \xrightarrow{\partial_1} C_0\) and
\(C^0 \xrightarrow{d_0} C^1 \xrightarrow{d_1} C^2\) over \(\mathbb{Q}\)
and \(\mathbb{F}_2\). \\
\textbf{Coboundary maps} & \(d_0\): gauge shift and boundary
identification; \(d_1\): topological cycle and loop closure detector. \\
\textbf{ZeroRel} & Invariance of cluster Boltzmann weights and
observable traces under spin and bond configurations. \\
\textbf{Licensed repairs} & Topological cluster involution \(\tau\)
reversing orientation of odd-winding meron clusters or string loops with
\((1 + (-1)) w(C) = 0\). \\
\textbf{Quotient / Q1} & Homology quotient
\(H_1 = \ker(\partial_1)/\mathrm{im}(\partial_2)\) classifying
non-contractible meron loops modulo boundaries. \\
\textbf{Residuals} & \(\Phi_1\): topological cycle sign obstruction;
\(\Gamma_2\): higher boundary discrepancy; duality gap \(\Delta = 0\)
over rational arithmetic. \\
\textbf{Witnesses} & Fixed-point-free sign-reversing involution
certificates \(\Xi\), GIN inverted index doublet pairings, and Lean 4
mechanized cancellation lemmas. \\
\textbf{Observers / profiles} & Exact cluster sum oracle,
transfer-matrix traces, and lexical hygiene validator
(\passthrough{\lstinline!assert\_lexical\_hygiene!}). \\
\textbf{Anomalons} & Single-meron and odd-winding topological defects
carrying non-trivial holonomy / Dirac sign phases. \\
\textbf{CompletionLaw} & Exact rational arithmetic over \(\mathbb{Q}\)
and finite group completion under \(\mathbb{Z}_2\),
\(\mathbb{Z}_2 \times \mathbb{Z}_2\), and \((\mathbb{Z}_2)^3\). \\
\textbf{Terminal / route} & Deterministic routing:
\passthrough{\lstinline!ExactPairwiseCancellation!},
\passthrough{\lstinline!PositiveSectorPreservation!},
\passthrough{\lstinline!ZeroDualityGapRepair!}, or
\passthrough{\lstinline!TypedRefusal!}. \\
\textbf{AuthorityEffect} & None. All results are finite computational or
machine-checked mathematical demonstrations; no ungrounded authority is
claimed. \\
\end{longtable}

\subsubsection{Summary of mechanisms and
evidence}\label{summary-of-mechanisms-and-evidence}

\begin{longtable}[]{@{}
  >{\raggedright\arraybackslash}p{(\linewidth - 2\tabcolsep) * \real{0.5000}}
  >{\raggedright\arraybackslash}p{(\linewidth - 2\tabcolsep) * \real{0.5000}}@{}}
\toprule\noalign{}
\begin{minipage}[b]{\linewidth}\raggedright
Mechanism
\end{minipage} & \begin{minipage}[b]{\linewidth}\raggedright
Status and scope
\end{minipage} \\
\midrule\noalign{}
\endhead
\bottomrule\noalign{}
\endlastfoot
Finite orbit factorization and cancellation & Conditional algebra under
Assumptions 2.1 and 2.2 \\
Per-orbit sign identities & Lean statements listed in Section 7; no
SQL-to-Lean bridge \\
SQLite construction and independent transfer trace & Exact computation
for the Section 3 fixture (\(Z = 912\)) and the Section 6.2 six-site run
(\(Z = 33{,}012\)) \\
Historical prompt filter and scorer & Tested program behavior; retained
pilot has answer leakage \\
New finite witness protocol & Four selected orbits, complete supplied
inputs and checked pair witnesses \\
Multi-stage DDC DAG and Gamma38 wire grammar & Verified 5-stage causal
DAG in JSONL and Gamma38 wire packets with zero domain keywords (Section
6.4) \\
Simplicial boundary interpretation & Unfulfilled typed transport
obligation (Section 8.1) \\
2D staggered Gross-Neveu meron cluster solver & Exact rational
cancellation \((1 + (-1)) w(C) = 0\) with zero duality gap on \(T^2\)
(\(4 \times 4\)), \(Z = 150\) (Section 6.5) \\
3D closed meron string loops and Gauss linking & Exact rational
cancellation with zero duality gap on \(T^3\) (\(6 \times 6 \times 6\)),
\(Z = 310\) (Section 6.6) \\
OpenRouter 2D meron cluster model benchmark & Multi-turn wire evaluation
across 3 model tiers in zero-shot/few-shot modes (Section 6.7) \\
Multi-turn interactive PFP trials & Coached independent prompt runner
and synthetic model-labelled fixture; provider measurements unverified
(Section 6.4) \\
Pure \(\Gamma_{38}\) symbolic wire runner & Finite host replay and exact
expected-packet checks; Wire-only input; offline candidates are supplied
fixtures (Section 6.4) \\
\end{longtable}

\subsection{9. Code and data}\label{code-and-data}

Paths are relative to the package root, which is
\nolinkurl{omega_engine_projects/fsp/meron_sign_problem/} in the source
repository
\href{https://github.com/dbccpage/the_symbol_game}{dbccpage/the\_symbol\_game}.

\begin{itemize}
\tightlist
\item
  \nolinkurl{code/meron_exact_model.py} and
  \passthrough{\lstinline!code/meron\_schema.sql!}: the exact fixture,
  copied from
  \nolinkurl{omega_engine_projects/fsp/code/sign_sql/exact_models.py}
  and \nolinkurl{omega_engine_projects/fsp/databases/meron_hubbard.sql}
  at commit \nolinkurl{52cb5457f6c485ba548cfd292f6012a941ed22cf.} Only
  the meron functions and tables were copied; the file headers list the
  changes.
\item
  \passthrough{\lstinline!code/oracle.py!},
  \passthrough{\lstinline!code/blind\_problem.py!},
  \passthrough{\lstinline!code/prompt\_builder.py!},
  \passthrough{\lstinline!code/antibody\_miner.py!} and
  \nolinkurl{code/run_blind_benchmark.py}: the oracle wrapper, the
  filter and problem text, the prompt builders, the scorer and the
  harness.
\item
  \nolinkurl{code/witness_benchmark.py} and
  \nolinkurl{code/run_witness_benchmark.py}: the separate finite witness
  tasks and bounded offline/live runner described in Section 6.1.
\item
  \nolinkurl{docs/witness_benchmark.md}: protocol, scope, and cochain
  transport notes.
\item
  \passthrough{\lstinline!code/results/!}: the retained run record, the
  anonymous model's output, and certified PFP ladder and scaling
  receipts.
\item
  \passthrough{\lstinline!code/tests/!}: the tests cited in the claim
  ledger.
\item
  \passthrough{\lstinline!code/symbol\_game\_pfp.py!}: the canonical
  Symbol Game PFP implementation.
\item
  \passthrough{\lstinline!demos/!}: runnable replay scripts
  (\nolinkurl{demo_4site_symbol_game.py},
  \nolinkurl{demo_6site_scaling_gin.py},
  \nolinkurl{demo_multistage_symbol_game.py},
  \nolinkurl{demo_multiturn_interactive_pfp.py},
  \nolinkurl{demo_pure_symbolic_meron.py},
  \nolinkurl{demo_pure_symbolic_6site.py},
  \nolinkurl{demo_2d_gross_neveu_meron.py},
  \nolinkurl{demo_3d_meron_string_linking.py},
  \nolinkurl{bench_openrouter_meron_2d.py}), demo test suite
  (\passthrough{\lstinline!demos/tests/!}), and certified receipts
  (\passthrough{\lstinline!demos/results/!}).
\item
  \passthrough{\lstinline!lean4/!}: the Lean 4 project.
\end{itemize}

The run record was produced by the harness in the repository directory
\nolinkurl{omega_engine_projects/fsp/demos/blind_cluster_recovery/} and
committed in e2a1e9fd0. Commit 8c0a0d40d then changed the scorer
(Section 5.1), and the package copy followed in 705f9cbb1. The prompt
strings have not changed since the run.

To reproduce the checks, run from the package root with Python 3.12:

\begin{lstlisting}
py -3.12 -m pytest code/tests/test_meron_cluster_solution.py -q
py -3.12 -m pytest code/tests/test_pilot_evidence.py -q
py -3.12 -m pytest code/tests/test_space_bunny.py -q
py -3.12 -m pytest demos/tests/test_demo_2d_gross_neveu.py -q
py -3.12 -m pytest demos/tests/test_demo_3d_meron_string.py -q
py -3.12 -m pytest demos/tests/test_bench_openrouter_meron_2d.py -q
py -3.12 -m ruff check code demos/code demos/tests
py -3.12 code/run_blind_benchmark.py --dry-run --output-dir <new-directory>
cd lean4 && lake build && lake env lean Census.lean
\end{lstlisting}

The dry run replaces every model call with a fixed answer and makes no
network request (Claim MSP-16); the budget guard blocks requests
exceeding the configured cost ceiling (Claim MSP-16); choose an output
directory outside the package. A live run needs authorization and an
OpenRouter key and is not part of these checks.

\subsection{Verification for this
revision}\label{verification-for-this-revision}

The package verifier (\passthrough{\lstinline!verify\_deposit.py!})
checks the manifest and the two-way file census, and
\passthrough{\lstinline!tools/rebuild\_release.py --check!} checks the
release archive locally (Claim MSP-33). The package tests, lint, and
Lean build are executed by the dated review receipt recorded with this
revision; that receipt lists each command, tool version, exit status,
and tested scope. Historical receipts in this package keep their
original dates and scope and are not current execution evidence.

\subsection{10. Claim index}\label{claim-index}

The following table indexes all claims MSP-01 through MSP-37 declared in
\passthrough{\lstinline!paper\_claim\_ledger.md!}, reconciling asserted
propositions against their mechanized proofs, exact computational
oracles, and empirical receipts.

\begin{longtable}[]{@{}
  >{\raggedright\arraybackslash}p{(\linewidth - 6\tabcolsep) * \real{0.2500}}
  >{\raggedright\arraybackslash}p{(\linewidth - 6\tabcolsep) * \real{0.2500}}
  >{\raggedright\arraybackslash}p{(\linewidth - 6\tabcolsep) * \real{0.2500}}
  >{\raggedright\arraybackslash}p{(\linewidth - 6\tabcolsep) * \real{0.2500}}@{}}
\toprule\noalign{}
\begin{minipage}[b]{\linewidth}\raggedright
ID
\end{minipage} & \begin{minipage}[b]{\linewidth}\raggedright
Claim
\end{minipage} & \begin{minipage}[b]{\linewidth}\raggedright
Status
\end{minipage} & \begin{minipage}[b]{\linewidth}\raggedright
Evidence
\end{minipage} \\
\midrule\noalign{}
\endhead
\bottomrule\noalign{}
\endlastfoot
MSP-01 & For one orbit with unit weight, if some cluster is a meron the
signs sum to 0 (Lemma 2.4). & Lean theorem &
\nolinkurl{MeronCharacter.meron_orbit_cancels}; pairwise
\nolinkurl{MeronClusterInvolution.meron_pair_exact_cancellation} \\
MSP-02 & For one orbit with unit weight and no meron, the signs sum to
\(2^k\) (Lemma 2.5). & Lean theorem &
\nolinkurl{MeronCharacter.zero_meron_orbit_positive} \\
MSP-03 & One orbit sums to 0 if it has a meron and to \(2^k\) otherwise
(Theorem 2.6 per orbit). The factor \(W_0(g)\) and sum over graphs are
not mechanized. & Lean theorem, scoped &
\nolinkurl{MeronCharacter.compression_preserves_partition} \\
MSP-04 & A weight-preserving, sign-reversing involution has no fixed
point, whatever the weight. & Lean theorem &
\nolinkurl{MeronClusterInvolution.no_fixed_points_when_sign_reverses},
\nolinkurl{fixed_point_contradiction} \\
MSP-05 & The fixture (4 sites, 2 repetitions, \(q=3\)) has 256 graphs,
1,296 configurations, 192 negative, 57 cancelling graphs, 199 graphs
without meron, \(Z=912\), \(Z_{\mathrm{abs}}=1,296\), average sign
19/27. Orbits of the 57 sum to 0; 199 sum to 912. & Exact computation &
\nolinkurl{test_meron_cluster_solution.py::test_oracle_ground_truth_invariants} \\
MSP-06 & An occupation-basis transfer-matrix trace, computed without
clusters, also gives 912. & Exact computation &
\nolinkurl{test_meron_cluster_solution.py::test_cluster_sum_matches_independent_transfer_trace} \\
MSP-07 & On the fixture, Assumptions 2.1 and 2.2 hold: weights are
constant on orbits and sign is product of single-flip signs. & Exact
computation &
\nolinkurl{test_meron_cluster_solution.py::test_fixture_satisfies_character_law_and_equal_orbit_weights} \\
MSP-08 & The filter rejects text containing any of twelve substrings;
all five prompt builders pass it. Prompts keep other sign-problem
vocabulary. & Tested behavior &
\nolinkurl{test_meron_cluster_solution.py::test_lexical_hygiene_blocks_all_forbidden_tokens} \\
MSP-09 & Retained record lists 22 calls to 5 models with Table 2 scores;
SHA-256 digest in \passthrough{\lstinline!receipt.json!}; estimated
spend 0.01496 USD is unbilled. & Retained record &
\nolinkurl{test_pilot_evidence.py::test_recorded_call_ledger};
\nolinkurl{test_pilot_evidence.py::test_receipt_digest_matches_report} \\
MSP-10 & Re-scoring 12 retained single-turn answers with package scorer
reproduces every score. Rubric checks values and keywords, not
derivations. & Tested behavior &
\nolinkurl{test_pilot_evidence.py::test_recorded_scores_reproduce_under_package_scorer} \\
MSP-11 & No cell of Arms 1 to 3 passed (0 of 9). Prompts state 1,296 and
19/27 fixing \(Z=912\) and state the target property; 4 samples have no
odd-winding cycle. & Retained record &
\nolinkurl{test_pilot_evidence.py::test_unassisted_prompts_state_quantities_that_fix_z} \\
MSP-12 & GLM 5.3 Flash in Arm 2 wrote \(Z=912\) under key
\passthrough{\lstinline!Z!} outside the schema and scored 0. & Retained
record &
\nolinkurl{test_pilot_evidence.py::test_glm_arm2_answer_used_keys_outside_the_schema} \\
MSP-13 & Arm 4 feedback states ``Z must equal 912'' and the 57/199
split; Arm 5 prompt states sign formula and cancellation step. Every
passing run had such a prompt. & Tested prompts &
\nolinkurl{test_pilot_evidence.py::test_arm4_corrections_state_the_target} \\
MSP-14 & Arm 3 prompt has 1,124 characters vs 4,320 for Arm 1 by
omitting narrative and configuration lists. & Tested prompts &
\nolinkurl{test_pilot_evidence.py::test_arm3_prompt_is_shorter_because_it_omits_content} \\
MSP-15 & Anonymous model output records no model, prompt, parameters,
latency or score; re-scores 4/4 outside schema; counts 57 and 199 come
from corrections. Not reproducible. & Retained record &
\nolinkurl{test_space_bunny.py::test_retained_record_fields} \\
MSP-16 & Harness dry run makes no network call; budget guard blocks
requests over configured price ceiling. & Tested behavior &
\nolinkurl{test_meron_cluster_solution.py::test_dry_run_executes_safely_without_spend} \\
MSP-17 & Language models discovered cancellation mechanism, or prompt
formats enabled discovery. & Withdrawn & Contradicted by MSP-11 and
MSP-13 \\
MSP-18 & Prompt compression reduced latency to 1.6 s. & Withdrawn &
Llama 3.3 70B Arm 3 scored 0 after 133.7 s; latency not attributable to
prompt \\
MSP-19 & Four selected orbits at \(q=5/2\) form a reproducible witness
task with complete input data, opaque IDs, no supplied answers, and
exact checks. & Finite tested task &
\nolinkurl{code/witness_benchmark.py};
\nolinkurl{docs/witness_benchmark.md} \\
MSP-20 & Host positive control succeeds with zero provider calls. Not a
model trial. & Offline execution &
\nolinkurl{code/run_witness_benchmark.py}; host receipt \\
MSP-21 & Cocycle evaluates equally on cycles differing by a boundary.
Does not construct Meron-to-cochain map. & Conditional algebra & Section
8.1; Hatcher Section 3.1 \\
MSP-22 & Retained live-trial results claims withdrawn; only committed
receipt is host positive control with 0 provider calls. & Withdrawn &
\nolinkurl{docs/witness_host_control_2026_10_02/results.json} \\
MSP-23 & Involution mechanization claims withdrawn as stated;
\nolinkurl{MeronClusterInvolution.lean} covers abstract pair only, not
coordinate synthesis or URN packets. & Withdrawn &
\nolinkurl{lean4/MeronClusterInvolution.lean} \\
MSP-24 & \passthrough{\lstinline!meron\_database(6, 2, "3")!} returns
4,096 graphs, 33,396 configurations, 192 negative, 4,024 zero-meron
graphs (72 with positive merons), partition 33,012, transfer trace
33,012. & Exact computation &
\nolinkurl{test_meron_cluster_solution.py::test_six_site_q3_census_matches_transfer_trace} \\
MSP-25 & Under Symbol Game PFP with DDC cycle spectra and host antibody
guards, 4 models score 4/4 against BlindClusterOracle on 4-site fixture.
& Certified empirical trial &
\nolinkurl{demos/code/demo_4site_symbol_game.py};
\nolinkurl{test_symbol_game_pfp_ladder.py} \\
MSP-26 & On 6-site fixture, blind DDC multi-meron loop identification
detects 72 cancelling subgraphs and 4,024 surviving subgraphs; GIN index
pairs all 384 cancelling states into 192 doublets, \(Z=33,012\). & Exact
computation & \nolinkurl{demos/code/demo_6site_scaling_gin.py};
\nolinkurl{test_symbol_game_6site_pfp.py} \\
MSP-27 & DDC 5-stage pipeline emits append-only DAG in JSON Lines with
\passthrough{\lstinline!LIFT\_REQUEST!}/\passthrough{\lstinline!LIFT\_GRANT!}
and \(\Gamma_{38}\) wire encoding, reproducing exact repair \(Z=912\) (4
sites) and \(Z=33,012\) (6 sites) under lexical hygiene. & Protocol
verification & \nolinkurl{code/symbol_game_stage_controller.py};
\nolinkurl{demo_multistage_symbol_game.py} \\
MSP-28 & Physical sign-factorization and involution witnesses
demonstrated on finite 2D/3D lattices; continuum scaling and doped
regimes open. & Tested finite models & Section 8.2;
\nolinkurl{demo_2d_gross_neveu_meron.py};
\nolinkurl{demo_3d_meron_string_linking.py} \\
MSP-29 & Coached prose runner retains raw responses; retained multi-turn
record is synthetic fixture; provider latency and discovery rates
unverified. & Tested runner &
\nolinkurl{demos/code/demo_multiturn_interactive_pfp.py} \\
MSP-30 & Deterministic host recomputes \(Z=912\) and transfer trace;
revised 4-site inputs are wire-only; offline responses supplied fixtures
with LlmTurn bindings. & Host computation &
\nolinkurl{demos/code/demo_pure_symbolic_meron.py} \\
MSP-31 & Retained 3-model summary contains 4 expected wire examples per
model but lacks raw responses; latency, usage, and discovery unverified.
& Retained summary &
\nolinkurl{demos/results/certified_pure_wire_ladder_results.json} \\
MSP-32 & On 6-site fixture, host enumerates 4,096 graphs, 33,396
configurations, pairs 384 cancelling states into 192 doublets, partition
and transfer trace 33,012. & Host computation &
\nolinkurl{demos/code/demo_pure_symbolic_6site.py} \\
MSP-33 & Package manifests and deterministic ZIP checked locally;
retained phase signoffs are authored assertions with authority effect
None. & Local inventory checks &
\passthrough{\lstinline!verify\_deposit.py!};
\passthrough{\lstinline!tools/rebuild\_release.py!} \\
MSP-34 & For repetitions 2 and \(q=3\), bridge stores \(Z(4)=912\),
\(Z(6)=33,012\), \(Z(8)=973,824\); Neville extrapolation midpoint agrees
with 1.69411 through 4 digits. & Finite traces &
\nolinkurl{omega_engine_projects/.../meron_wilsonian_bridge.py}; Section
6.8 \\
MSP-35 & On 2D staggered Gross-Neveu on \(T^2\) (8 sites, 4 bond graphs,
16 states), cluster breakup finds 2 odd-winding meron graphs whose 8
states cancel pairwise under \(\tau\), yielding \(Z_{\mathrm{odd}}=0\)
and \(Z=150\), matching transfer trace 150. & Exact computation &
\nolinkurl{demos/code/demo_2d_gross_neveu_meron.py}; Section 6.5 \\
MSP-36 & On 3D periodic \(T^3\) (8 sites, 12 links), closed meron loops
classified by winding parity \(\mathbf{W} \in (\mathbb{Z}_2)^3\) and
Gauss linking \(Lk \in \mathbb{Z}\); flip involution cancels all 4
odd-winding configurations pairwise to 0, yielding \(Z=310\). & Exact
computation & \nolinkurl{demos/code/demo_3d_meron_string_linking.py};
Section 6.6 \\
MSP-37 & OpenRouter multi-turn model ladder benchmark on 2D meron
cluster breakup evaluates frontier models across 4 stages in pure
\(\Gamma_{38}\) wire format under zero-smuggling lexical hygiene. &
Empirical benchmark &
\nolinkurl{demos/code/bench_openrouter_meron_2d.py}; Section 6.7 \\
\end{longtable}

\subsection{References}\label{references}

\begin{enumerate}
\def\labelenumi{\arabic{enumi}.}
\tightlist
\item
  S. Chandrasekharan and U.-J. Wiese, ``Meron-Cluster Solution of
  Fermion Sign Problems,'' Phys. Rev.~Lett. \textbf{83}, 3116 (1999).
  DOI:
  \href{https://doi.org/10.1103/PhysRevLett.83.3116}{\nolinkurl{10.1103/PhysRevLett.83.3116}}.
  arXiv:
  \href{https://arxiv.org/abs/cond-mat/9902128}{cond-mat/9902128}.
\item
  M. Troyer and U.-J. Wiese, ``Computational Complexity and Fundamental
  Limitations to Fermionic Quantum Monte Carlo Simulations,'' Phys.
  Rev.~Lett. \textbf{94}, 170201 (2005). DOI:
  \href{https://doi.org/10.1103/PhysRevLett.94.170201}{\nolinkurl{10.1103/PhysRevLett.94.170201}}.
\item
  E. Huffman and S. Chandrasekharan, ``Solution to sign problems in
  half-filled spin-polarized electronic systems,'' Phys. Rev.~B
  \textbf{89}, 111101 (2014). DOI:
  \href{https://doi.org/10.1103/PhysRevB.89.111101}{\nolinkurl{10.1103/PhysRevB.89.111101}}.
  arXiv: \href{https://arxiv.org/abs/1311.0034}{1311.0034}.
\end{enumerate}

\end{document}
